\section{The fixed cubic Fourier packet and the actual local normalization}
\label{sec:cubic-local-fourier}

The obstruction in AC-G40 is a literal hypothesis obstruction, not a
failure of finite Fourier analysis at the primes above $269$.  The
following calculation determines the relevant cubic matrix Gauss
factor exactly, including the denominator in the source normalization.

\begin{theorem}[AC-G41: the cubic matrix Fourier packet]
\label{AC-G41}
Let $p\equiv-1\pmod3$ be prime, let $f>0$ be even, put
$k=\mathbb F_{p^f}$ and $q=p^f$, and let
$\chi:k^\times\longrightarrow\mu_3$ be nontrivial.  Extend
$\chi$ and $\chi^{-1}$ by zero at zero.  Set
\[
 \psi(x)=\exp\!\left(\frac{2\pi i}{p}
                 \operatorname{Tr}_{k/\mathbb F_p}(x)\right),
 \qquad
 G(\chi)=\sum_{x\in k}\chi(x)\psi(x).
\]
There is a sign $\epsilon\in\{1,-1\}$ such that
\[
 G(\chi)=G(\chi^{-1})=\epsilon p^{f/2}.                 \tag{G41.1}
\]
For $r\geq1$, put $\Phi_\chi(X)=\chi(\det X)$ on
$M_r(k)$ and define, for $a\in k^\times$,
\[
 (\mathcal F_a u)(Y)=q^{-r^2/2}
          \sum_{X\in M_r(k)}u(X)\psi(a\operatorname{tr}(XY)).
\]
Then
\[
 \mathcal F_a\Phi_\chi
       =\epsilon^r\chi(a)^{-r}\Phi_{\chi^{-1}},
 \qquad
 \mathcal F_a\Phi_{\chi^{-1}}
       =\epsilon^r\chi(a)^r\Phi_\chi.                  \tag{G41.2}
\]
In fact, on their two-dimensional complex span $W$,
\[
       \mathcal F_a=\epsilon^r P_{a,p},\qquad
       (P_{a,p}u)(Y)=u((aY)^{(p)}),                    \tag{G41.3}
\]
where the superscript denotes entrywise $p$-th power.
Consequently both
\[
 R\Phi_\chi+R\Phi_{\chi^{-1}}
 \quad\hbox{and}\quad W\cap R^{M_r(k)}
                                                               \tag{G41.4}
\]
are preserved by $\mathcal F_a$ and its inverse for every subring
$R\subset\mathbb C$ containing $\mathbb Z[\omega]$.
This applies to the localizations of rings of integers at every
coefficient prime, including $3$ and $p$, and hence after their
completion.  No division by $3$ or assertion that the first lattice
in (G41.4) is saturated is used.
\end{theorem}
\begin{proof}
\proofstep{1. The scalar sum.}
The trace is unchanged by $x\mapsto x^p$, and
$\chi(x^p)=\chi(x)^{-1}$.  Changing variables gives
$G(\chi)=G(\chi^{-1})$.  Additive orthogonality gives
\begin{align*}
 G(\chi)G(\chi^{-1})
  &=\sum_{t\in k^\times}\chi(t)
                  \sum_{y\in k^\times}\psi((t+1)y)\\
  &=(q-1)\chi(-1)
                  -\sum_{t\in k^\times\setminus\{-1\}}\chi(t)
    =q\chi(-1)=q.
\end{align*}
Here the sum of a nontrivial multiplicative character is zero,
and $\chi(-1)=1$ because its order divides both $2$ and $3$.
Thus the common Gauss sum squares to $q$, proving (G41.1).
Replacing $\psi(x)$ by $\psi(ax)$ multiplies $G(\chi)$ by
$\chi(a)^{-1}$.

\proofstep{2. The matrix sum, without an external formula.}
Write
\[
       G_r(\chi)=\sum_{X\in M_r(k)}
                          \chi(\det X)\psi(\operatorname{tr}X).
\]
Fix the first $r-1$ rows and let $c\in k^r$ be their cofactor
vector, so that the determinant with last row $z$ is $z\cdot c$.
If $c=0$, every contribution is zero.  If $c$ is not a nonzero
multiple of the last coordinate vector, there is a vector
$v\in\ker(c\cdot-)$ with $v_r\ne0$.  A scalar multiple of $v$
has $\psi(v_r)\ne1$.  Translation of the last row by this vector
preserves the determinant and multiplies the summand by
$\psi(v_r)$, so the sum over last rows vanishes.

In the remaining case the first $r-1$ rows are $(B\mid0)$ with
$B\in\mathrm{GL}_{r-1}(k)$ and $c_r=\det B$.  Their contribution
after summing over the last row is
\[
 q^{r-1}\chi(\det B)\psi(\operatorname{tr}B)G(\chi).
\]
Induction, starting with $r=1$, proves
\[
 G_r(\chi)=q^{r(r-1)/2}G(\chi)^r
                       =\epsilon^r q^{r^2/2}.          \tag{G41.5}
\]
The same formula holds for $\chi^{-1}$, with the same sign.

\proofstep{3. Every Fourier value.}
If $Y$ is invertible, the substitution $Z=aXY$ gives
\[
 \sum_X\chi(\det X)\psi(a\operatorname{tr}(XY))
       =\chi(a)^{-r}\chi(\det Y)^{-1}G_r(\chi).
\]
If $Y$ is singular, choose $A\in\mathrm{GL}_r(k)$ fixing
$\operatorname{im}Y$ pointwise and having $\chi(\det A)\ne1$.
Such an $A$ exists by prescribing a nontrivial determinant on
a complementary subspace.  The substitution $X\mapsto XA$
preserves $\operatorname{tr}(XY)$ and multiplies the determinant
character by $\chi(\det A)$; the sum is therefore zero.
This proves (G41.2), including its values on singular matrices.

Since $p\equiv-1\pmod3$,
\[
 \Phi_\chi((aY)^{(p)})
                  =\chi(a)^{-r}\Phi_{\chi^{-1}}(Y),
\]
and the analogous identity holds with $\chi^{-1}$.  The two
functions are linearly independent: determinant surjectivity
reduces that assertion to the independence of two distinct
characters on $k^\times$.  Hence (G41.3) follows.  The map
$Y\mapsto(aY)^{(p)}$ permutes the finite coordinate set and
preserves $W$.  It preserves both lattices in (G41.4), as does
its inverse.  This proves all the integrality assertions.
\end{proof}

\begin{corollary}[AC-G41-normalization: the source factor at every coefficient prime]
\label{AC-G41-normalization}
For $0\leq r\leq n$, use the usual empty-matrix convention if
$r=0$.  The precise factor appearing in the source's
\path{GlobalTheory.GaussUnit.unitOff_prod_gaussSum} is
\[
 C_{n,r}(k,\chi,a)
   :=q^{-nr}\sum_{g\in\mathrm{GL}_r(k)}
                  \chi(\det g)^{-1}\psi(a\operatorname{tr}g)
   =\epsilon^r\chi(a)^r p^{-f r(2n-r)/2}.              \tag{G41.6}
\]
Thus its normalized additive valuation is zero at every
coefficient prime $\ell\ne p$, and at $\ell=p$ is exactly
$-fr(2n-r)/2$.  Its absolute logarithmic height is
$fr(2n-r)\log(p)/2$.

Let $F\supset\mathbb Q(\omega)$ be unramified at $269$, let
$Q$ be all primes of $F$ above $269$, and use nontrivial cubic
residue characters at those primes.  With fixed $n,r$, the
product of the factors (G41.6) is
\[
                \zeta\,269^{-B},\qquad
 B=\frac{[F:\mathbb Q]r(2n-r)}2,
                   \qquad \zeta\in\mu_6.              \tag{G41.7}
\]
In particular this entire product is a $3$-adic unit.  Its
source normalization is not asserted to be a unit at $269$.
\end{corollary}
\begin{proof}
Apply (G41.5) to $\chi^{-1}$ and use the scaling formula in
step 1.  This gives (G41.6).  The exponent is an integer
because $f$ is even; all remaining factors are roots of unity.
Valuations and heights follow.  For (G41.7), AC-G40 proves
that every residue degree is even.  Unramifiedness gives
$\sum_{P\mid269}f(P/269)=[F:\mathbb Q]$.  Multiply (G41.6)
and add the exponents.
\end{proof}

\begin{proposition}[AC-G41-trace: the unramified trace replacement for prime-field digits]
\label{AC-G41-trace}
Let $L/\mathbb Q_p$ be unramified of degree $f$, with valuation
ring $\mathcal O_L$.  For $D\geq1$ put
$A_D=\mathcal O_L/p^D\mathcal O_L$.  There is an integral
$\mathbb Z_p$-basis $b_1,\ldots,b_f$ with an integral trace-dual
basis $b_1^*,\ldots,b_f^*$.  In particular
\[
 A_D\simeq(\mathbb Z/p^D\mathbb Z)^f,
 \quad
 (x,y)\longmapsto
  \exp\!\left(\frac{2\pi i}{p^D}
              \operatorname{Tr}_{L/\mathbb Q_p}(xy)\right)
                                                               \tag{G41.8}
\]
is a perfect pairing of finite additive groups.  Multiplication
by any unit in its trace argument preserves perfectness.

On $A_D^s$ the corresponding Fourier transform with factor
$p^{-fDs/2}$ is unitary and its square is the reflection
$u(x)\mapsto u(-x)$.  If $f$ is even, its matrix and inverse
are integral at every coefficient prime $\ell\ne p$ after
adjoining the $p^D$-th roots of unity.  Thus the finite additive
Fourier and residue-coordinate construction needed at an
unramified prime has a direct replacement for arbitrary $f$.

When $f>1$, the stronger assertion used by the source's
prime-field digit construction, that every integral local
element has a rational integer representative modulo $p$,
is false.  Also $q=p^f$ is not a uniformizer:
$q\mathcal O_L=p^f\mathcal O_L\subset p^2\mathcal O_L$.
Consequently this proposition changes the local interface;
it does not prove the literal predicate \texttt{TwDegOne}.
\end{proposition}
\begin{proof}
Choose lifts of a residue-field basis over $\mathbb F_p$.
They form a $\mathbb Z_p$-basis of $\mathcal O_L$: their
span surjects modulo $p$, hence equals $\mathcal O_L$ by
Nakayama, and there are exactly $f$ generators in a free
module of rank $f$.  The trace Gram matrix reduces to the
nondegenerate trace pairing of $\mathbb F_{p^f}/\mathbb F_p$.
Indeed its trace polynomial
$X+X^p+\cdots+X^{p^{f-1}}$ is nonzero and has degree less
than $p^f$, so it cannot vanish on the entire field.  For
any nonzero $x$, multiplication by $x$ is a bijection;
some $y$ therefore has nonzero trace of $xy$.
The Gram determinant is a $p$-adic unit.
The inverse Gram matrix therefore has
entries in $\mathbb Z_p$, giving the integral dual basis.

In these two bases the pairing in (G41.8) is the ordinary
dot-product pairing on $(\mathbb Z/p^D\mathbb Z)^f$.
It is perfect, and a unit in the argument is an automorphism
of $A_D$.  Orthogonality of its characters gives, on $A_D^s$,
\[
 \sum_{y\in A_D^s}
       \exp\!\left(\frac{2\pi i}{p^D}
           \operatorname{Tr}((x+z)\cdot y)\right)
 =\begin{cases}p^{fDs},&z=-x,\\0,&z\ne-x.\end{cases}
\]
This proves the square formula and, with complex conjugation,
unitarity.  For even $f$, the normalizing factor is an integral
power of $p^{-1}$; the other entries are roots of unity.
The same holds for the inverse, proving integrality away from
$p$.  Finally choose a residue class outside the prime field
and lift it to $\mathcal O_L$.  No rational integer represents
it modulo $p$, proving the stated failure.  The assertion
about $q\mathcal O_L$ follows from unramifiedness and $f\geq2$.
\end{proof}

\paragraph{Exact source comparison and remaining gate.}
The pinned source at commit
\texttt{7cbf7a0c32fca1badf05931879a37c3f8a2b03b4},
\path{Development/AbelianVarietyTheory/Proofs/DivisorProofs/IsUnitAwayFromLevelProdNormalisedMatrixGaussSum.lean},
theorems \texttt{unitOff\_normalised} and
\texttt{unitOff\_prod\_gaussSum}, uses $q^{-nr}$ exactly as
in (G41.6).  Its general finite-field statement already allows
arbitrary residue degree and proves an away-from-level unit
property; (G41.6) is the more precise cubic evaluation, not
a correction to that statement.

In contrast,
\path{Formalism/LocalToGlobal/Columns/ChainTwOk.lean},
\texttt{TwDegOne}, requires both primality of
\texttt{qOf P} and \texttt{qOf P} outside $P^2$.
The same file's \texttt{TwCstQ}, and
\path{Formalism/LocalToGlobal/Columns/ChainTwPsi.lean},
\texttt{twRes\_exists}, \texttt{twPsiQ\_exists},
\texttt{twCstQ\_of}, and \texttt{twFin\_bdd}, use single
integer digits and powers of that same $q$.  The trace
construction above repairs the finite additive component
after distinguishing the uniformizer $p$ from the residue
cardinality $q$.  The remaining local norm, character,
and global column ties must still be checked for that
modified interface.

The source's
\path{Formalism/SharpRestriction/LevelLetterCocycles/IntegralHeckeForms.lean},
\texttt{IdxAt}, compares the complete projected theta
lattices uniformly after every finite DVR extension.
Neither (G41.2) nor the unit assertion at $3$ identifies
those complete lattices.  In particular the fixed cubic
\texttt{partD} term from AC-G40 is still present.  The
calculation removes a possible local Gauss-denominator
loss at coefficient $3$ and gives the exact loss at $269$;
it does not supply \texttt{GoodTw}, \texttt{NStage},
or the relative Selmer tie.

\paragraph{Representation dictionary.}
The cubic character is the finite rank-one representation
associated with the finite etale cover in the already audited
SGA~1, Expose~V, \S\S4 and~7 dictionary.  The determinant
functions here are explicit vectors in a finite function
space.  Their integral Fourier invariance neither asserts
the existence of a new motive nor changes the earlier
geometric and cohomological gates.
